
\section{Introduction}
\label{sec:introduction}

Soil characteristics and vegetation conditions are closely related
components of terrestrial ecosystems. Understanding their spatial
relationships is important for agricultural assessment, land-resource
management, and environmental monitoring. Geographic Information
Systems (GIS) and satellite remote sensing provide complementary
sources of information for investigating these relationships across
different spatial and temporal scales.

Conventional soil maps describe the spatial distribution of soil
types and associated characteristics. However, integrating these
datasets with contemporary geospatial information may require
coordinate reference system harmonization, geometry validation,
spatial clipping, attribute interpretation, and assessment of
mapping completeness. These operations are particularly relevant
when working with historical cartographic datasets that were
originally developed using formats and classification systems
different from those commonly employed in modern geospatial
workflows.

Satellite observations complement soil mapping by providing
repeated measurements of land-surface conditions. In particular,
the Sentinel-2 mission supplies multispectral imagery suitable
for vegetation monitoring. The Normalized Difference Vegetation
Index (NDVI), calculated from red and near-infrared reflectance,
provides an indirect indicator of vegetation greenness and canopy
development. Temporal NDVI observations can characterize changes
in vegetation conditions, although their interpretation depends
on land cover, acquisition quality, and agricultural context.

GeoSoil-Cuba is an open-source geospatial analysis project
developed to integrate soil mapping, land-cover information,
and Sentinel-2 satellite imagery within a reproducible
Python-based workflow. The project focuses on a study area
in Jaruco, Mayabeque Province, Cuba, using a 1:25,000-scale
soil map as the primary spatial framework for subsequent
soil--vegetation analysis.

The workflow has two complementary analytical components. The
first prepares and evaluates the soil dataset through geometry
validation and repair, coordinate reference system harmonization,
spatial intersection and clipping, classification interpretation,
and soil-map coverage assessment. Water-body information helps
distinguish areas excluded from terrestrial soil mapping from
unresolved gaps in the source dataset. The second component
processes Sentinel-2 observations to examine temporal NDVI
variability in a selected agricultural region and in three
locally verified potato-field samples. Together, these components
provide spatially organized datasets and exploratory vegetation
observations without presuming a causal relationship between
mapped soil classes and vegetation dynamics.

The main objectives of this study are to:

\begin{enumerate}
	
	\item Prepare and validate the existing soil cartographic
	information for geospatial analysis within the selected
	Area of Interest (AOI).
	
	\item Evaluate the spatial completeness and classification
	consistency of the soil dataset, distinguishing mapped
	water bodies from unresolved soil-data gaps.
	
	\item Integrate the processed soil information with
	land-cover classifications and organize it for analysis
	alongside Sentinel-2 observations.
	
	\item Process Sentinel-2 imagery and characterize temporal
	NDVI variability within a selected agricultural region.
	
	\item Compare NDVI observations from three locally verified
	potato-field samples, recognizing the limits of inferring
	crop development from satellite reflectance alone.
	
\end{enumerate}

The contribution of GeoSoil-Cuba lies primarily in the
integration, validation, and documentation of heterogeneous
geospatial datasets through an open-source analytical workflow.
Rather than proposing a new soil classification system or
establishing causal relationships between soil properties and
vegetation performance, the project provides a documented
framework for exploratory soil--vegetation analysis that can
be extended through additional satellite observations,
field measurements, and spatial datasets.
